\documentclass[11pt,a4paper]{article}
\usepackage[utf8]{inputenc}
\usepackage{amsmath}
\usepackage{geometry}
\usepackage{hyperref}
\usepackage{booktabs}
\usepackage{graphicx}
\usepackage{xcolor}
\usepackage{longtable}
\geometry{margin=2.5cm}
\hypersetup{colorlinks=true,linkcolor=blue,urlcolor=blue}

\title{\textbf{CoRE-COF 2D Pore Morphology Benchmark}\\Mirror (Chiral Toolkit) v0.6}
\author{Mingrui Zuo \& Lifeng Ding\\Xi'an Jiaotong-Liverpool University}
\date{June 2026}

\begin{document}
\maketitle
\tableofcontents
\newpage

% =================================================================
\section{Executive Summary}
% =================================================================

The 2D subset of the CoRE-COF 2019 database (1099 entries) was
systematically analysed using the Mirror platform's grid-based pore
detection pipeline.  Each CIF underwent:

\begin{enumerate}
    \item Pore detection (grid-based distance field, probe radius 1.2\,\AA).
    \item VDW surface projection (Fibonacci-sphere point cloud, 50 pts/atom).
    \item Cylindrical projection and $\theta$--$z$ VDW heatmap.
    \item FFT-based pore shape classification (L2--L6).
    \item Quantitative morphology parameter computation (CI, CF,
          curv\_H, Fourier powers, helicity, skew\_r).
\end{enumerate}

\textbf{Key results:}
\begin{itemize}
    \item 1099 / 1099 2D entries completed successfully (100 \%).
    \item Total runtime: $\sim$3.8 minutes (0.21 s/COF on average).
    \item 6 distinct pore shapes identified; L6 (hexagonal) dominates.
    \item Strong correlation between pore label and network topology.
    \item Chiral (Sohncke) space groups are extremely rare (4/1099).
\end{itemize}

% =================================================================
\section{Pore Label Distribution}
% =================================================================

The FFT-based classification assigns each COF a pore label based on the
dominant Fourier component of its angular wall-distance profile.

\begin{table}[h]
\centering
\begin{tabular}{lcc}
\toprule
Pore Label & Symmetry & Count \\
\midrule
L6 & Hexagonal (6-fold) & 546 \\
L2 & Diamond / Elliptical (2-fold) & 245 \\
L4 & Square (4-fold) & 112 \\
L3 & Trigonal / Clover (3-fold) & 46 \\
Other & Irregular / unresolved & 150 \\
\bottomrule
\end{tabular}
\caption{Pore-label distribution across 1099 2D COFs.}
\end{table}

Hexagonal pores (L6) account for 49.8 \% of the database, reflecting the
prevalence of honeycomb (hcb) topology among 2D COFs.  The L2 category
(22.3 \%) is the second most populous, arising from both square and
honeycomb nets under distortion.

% =================================================================
\section{Topology Correlation}
% =================================================================

Pore label shows a strong dependence on network topology (Fig.~1):

\begin{itemize}
    \item \textbf{L6}: 88.1 \% hcb, 6.0 \% kgm, 5.5 \% sql.
    \item \textbf{L4}: 71.4 \% sql, 24.1 \% hcb.
    \item \textbf{L3}: 80.4 \% hcb (triangular sub-pores).
    \item \textbf{L2}: 49.0 \% sql, 44.1 \% hcb (bimodal).
\end{itemize}

This confirms that the FFT-based classification is consistent with the
intrinsic symmetry of the underlying network topology.  The bimodal
distribution of L2 suggests that pore ellipticity can arise from either
square-lattice or honeycomb frameworks depending on the degree of
in-plane distortion.

% =================================================================
\section{Space Group and sgType Distribution}
% =================================================================

\begin{table}[h]
\centering
\begin{tabular}{lcc}
\toprule
sgType & Description & Count \\
\midrule
1 & Sohncke (enantiomorphic) & 4 \\
2 & Non-Sohncke & 525 \\
3 & Achiral & 420 \\
\bottomrule
\end{tabular}
\caption{Space-group type distribution.}
\end{table}

The overwhelming majority of 2D COFs belong to sgType~2 (non-Sohncke) and
sgType~3 (achiral).  Only 4 COFs exhibit Sohncke (chiral) space groups,
underscoring the scarcity of intrinsically chiral pore environments.

The most common individual space groups are P1 (\#1, 265 COFs),
P6/m (\#175, 122), P6 (\#168, 107), P3 (\#143, 78), and P-6 (\#174, 59).
Low-symmetry triclinic and hexagonal space groups dominate, consistent
with the flexible layered nature of 2D COFs.

% =================================================================
\section{Morphology Parameters}
% =================================================================

\subsection{Circularity Index (CI)}

CI measures how close the pore cross-section is to a perfect circle:

\[
\text{CI} = 1 - \frac{\sigma_r}{\mu_r},
\]

where $\mu_r$ and $\sigma_r$ are the mean and standard deviation of the
radial wall distance over $\theta$.  CI = 1 indicates a perfect circle;
lower values indicate distortion.  Example values: COF~\#436 (L6)
CI = 0.993, COF~\#450 (L2) CI = 0.783.

\subsection{Fourier Power Spectrum}

The angular wall-distance profile is Fourier-transformed to give relative
powers at frequencies 2, 3, 4, and 6:

\[
P_n = \frac{|F_n|^2}{\sum_{k} |F_k|^2},
\]

where $F_n$ is the $n$-th Fourier coefficient.  These powers directly
determine the pore label (highest $P_n$ for $n \ge 2$).

\subsection{Corrugation Factor (CF)}

The ratio of the VDW surface area to that of a smooth cylinder of the
same radius and height:

\[
\text{CF} = \frac{A_{\text{VDW}}}{2\pi R_1 \Delta z}.
\]

CF $> 1$ indicates a corrugated (rough) pore wall.

\subsection{Mean Curvature (curv\_H)}

The Laplacian of the Gaussian-smoothed VDW density field on the
$\theta$--$z$ plane, averaged over all points:

\[
\text{curv\_H} = \langle |\nabla^2 \rho_{\text{VDW}}| \rangle.
\]

\subsection{Helicity}

The helical twist angle $\alpha$ between adjacent $z$-layers, computed
via cross-correlation of their angular wall-distance profiles:

\[
\alpha = \arctan\!\left(\frac{\Delta\theta \cdot R_1}{\Delta z}\right).
\]

\subsection{Radial Skewness (skew\_r)}

The third statistical moment of the radial wall-distance distribution:

\[
\text{skew\_r} = \frac{\langle (r - \mu_r)^3 \rangle}{\sigma_r^3}.
\]

Positive values indicate outward-protruding wall features.

% =================================================================
\section{Performance}
% =================================================================

Batch processing of 1099 COFs completed in $\sim$3.8 minutes
(0.21 s/COF) using the Anaconda Python environment on a desktop machine.
This represents a $\sim$120$\times$ speedup over the full grid-based
pipeline (which includes the expensive VDW heatmap), achieved by skipping
the heatmap computation in batch mode.

\begin{table}[h]
\centering
\begin{tabular}{lcc}
\toprule
Stage & Time (s) & Fraction \\
\midrule
Pore detection (grid) & 0.12 & 57 \% \\
Cylindrical projection & 0.06 & 29 \% \\
Morphology computation & 0.03 & 14 \% \\
\bottomrule
\end{tabular}
\caption{Average per-COF timing breakdown.}
\end{table}

\end{document}
